\chapter{Cervical Cancer}

\section*{Definition and Overview}
Cervical cancer is a cancer that begins in the cells lining the cervix, the lower part of the uterus that connects to the vagina. Nearly all cases are caused by persistent infection with high-risk types of the human papillomavirus (HPV), which is spread through sexual contact.

Cervical cancer usually develops slowly over many years, passing through a stage of pre-cancerous cell changes that can be detected and treated before cancer forms. This makes it one of the most preventable and treatable cancers, thanks to HPV vaccination and regular cervical screening.

\section*{Epidemiology}
Cervical cancer is the fourth most common cancer in women worldwide, with the great majority of cases and deaths occurring in low- and middle-income countries where screening and vaccination are not widely available. It most often affects women in their 30s and 40s.

In countries with organised screening programs and HPV vaccination, incidence has fallen dramatically over recent decades. Australia was among the first countries to introduce the HPV vaccine, and cervical cancer is now projected to become rare in young vaccinated women.

\section*{Aetiology and Pathophysiology}
Almost all cervical cancers are caused by infection with high-risk HPV types, most commonly HPV 16 and HPV 18:

\begin{itemize}
    \item \textbf{HPV infection:} acquired through vaginal, anal, or oral sexual contact; very common, but usually cleared by the immune system within a year or two
    \item \textbf{Persistent infection:} when the immune system does not clear the virus, it can cause pre-cancerous change
    \item \textbf{Progression:} over years, low-grade changes may progress to high-grade changes and then to invasive cancer
    \item \textbf{Viral proteins:} HPV produces proteins that interfere with the cell's natural growth controls
    \item \textbf{Co-factors:} smoking, a weakened immune system, long-term contraceptive pill use, and having many sexual partners increase the risk
\end{itemize}

Squamous cell carcinoma accounts for most cases, with adenocarcinoma of the glandular cells making up most of the remainder.

\section*{Clinical Features}
Early cervical cancer and pre-cancerous changes often cause no symptoms at all, which is why screening is so important. When symptoms occur, they may include:

\begin{itemize}
    \item Abnormal vaginal bleeding: between periods, after sex (postcoital), or after the menopause
    \item Unusual vaginal discharge, which may be watery, pink, or foul-smelling
    \item Pain during intercourse
    \item Pelvic pain or pressure
    \item With advanced disease: pain in the back or legs, swelling of a leg, weight loss, fatigue, and problems passing urine or opening the bowels
\end{itemize}

\section*{Differential Diagnosis}
\begin{itemize}
    \item Cervical polyps, ectropion, or cervical erosion
    \item Cervicitis (inflammation or infection of the cervix), such as from chlamydia or gonorrhoea
    \item Vaginal infections, including thrush and bacterial vaginosis
    \item Fibroids and other benign uterine conditions
    \item Endometrial cancer, which also causes abnormal bleeding
    \item Bleeding related to pregnancy, contraception, or the perimenopause
\end{itemize}

\section*{When to Seek Medical Care}
See a doctor promptly if you have abnormal bleeding -- especially bleeding after sex, between periods, or after the menopause -- or an unusual vaginal discharge.

\textbf{Call 000 (or local emergency services) for:}
\begin{itemize}
    \item Heavy vaginal bleeding that is soaking through pads rapidly
    \item Severe pelvic or abdominal pain
    \item Fever with heavy bleeding or foul-smelling discharge
    \item Sudden collapse or signs of severe illness
\end{itemize}

\section*{Diagnosis and Investigations}
\begin{itemize}
    \item \textbf{Cervical screening test:} in Australia, a primary HPV test on a self- or clinician-collected swab every five years for women aged 25--74
    \item \textbf{Colposcopy:} a magnifying examination of the cervix using a special instrument
    \item \textbf{Cervical biopsy:} small samples of tissue are taken to confirm the diagnosis
    \item \textbf{HPV testing:} identifies high-risk types that warrant further investigation
    \item \textbf{Imaging:} MRI, CT, or PET scans to determine how far the cancer has spread (staging)
    \item \textbf{Blood tests:} including kidney and liver function as part of staging
\end{itemize}

\section*{Management}
Treatment depends on the stage of the cancer, the woman's age, and whether she wishes to preserve fertility:

\begin{itemize}
    \item \textbf{Pre-cancerous changes:} treatment by loop excision (LLETZ), laser, or freezing of the abnormal area
    \item \textbf{Early cancer:} cone biopsy, or surgery such as trachelectomy (removing the cervix while keeping the uterus) or hysterectomy
    \item \textbf{Radiotherapy:} external beam and internal (brachytherapy) radiation, often combined with chemotherapy
    \item \textbf{Chemotherapy:} used alongside radiation or for advanced disease
    \item \textbf{Targeted and immunotherapy:} newer drugs for advanced or recurrent disease
    \item \textbf{Multidisciplinary care:} a team of gynaecological oncologists, radiation oncologists, and specialist nurses plans treatment
    \item \textbf{Follow-up:} regular review, including examinations and imaging, to detect recurrence early
\end{itemize}

\section*{Complications}
\begin{itemize}
    \item Spread of cancer to nearby organs, the lymph nodes, or distant sites such as the lungs and liver
    \item Kidney damage or blockage of the ureters from tumour pressure
    \item Heavy bleeding from the tumour
    \item Fistula formation: an abnormal channel between the vagina and bladder or bowel
    \item Complications of treatment: menopausal symptoms from surgery or radiation, bladder and bowel irritation, lymphoedema of the legs, and infertility
    \item Psychological effects: anxiety, depression, and concerns about body image and sexuality
\end{itemize}

\section*{Prognosis and Outcomes}
Prognosis depends strongly on the stage at diagnosis. Pre-cancerous changes and very early cancers are usually cured, with excellent outcomes after simple treatment. Survival falls as the stage advances, but many women with locally advanced cancer are still treated successfully with combined radiation and chemotherapy. Screening and HPV vaccination are highly effective at preventing the disease altogether, and long-term survival after cervical cancer has improved steadily with better treatment and follow-up.

\section*{Prevention and Self-Care}
\begin{itemize}
    \item Have the HPV vaccine, ideally before sexual debut; it protects against the types that cause most cervical cancers
    \item Take part in cervical screening every five years from age 25 to 74, even if vaccinated
    \item Do not smoke, as smoking increases the risk of HPV persisting
    \item Use condoms to reduce the risk of HPV and other sexually transmitted infections
    \item Attend follow-up after any abnormal screening result or treatment
    \item Report abnormal bleeding or discharge promptly
    \item Support healthy immunity through a balanced diet, physical activity, and managing other health conditions
\end{itemize}

\section*{Sources and Further Reading}
The following organisations provide reliable information on cervical cancer, screening, and HPV.

\begin{itemize}
    \item National Health Service. \textit{NHS conditions guide on cervical cancer.} https://www.nhs.uk/conditions/
    \item Mayo Clinic. \textit{Cervical cancer: symptoms, causes, and treatment.} https://www.mayoclinic.org/
    \item World Health Organization. \textit{Global strategy to eliminate cervical cancer and HPV information.} https://www.who.int/
    \item Centers for Disease Control and Prevention. \textit{HPV and cervical cancer prevention information.} https://www.cdc.gov/
    \item MedlinePlus. \textit{Consumer health information on cervical cancer.} https://medlineplus.gov/
    \item MSD Manual. \textit{Cervical cancer -- professional overview.} https://www.msdmanuals.com/
\end{itemize}
